\chapter{Progettazione dell'applicazione}
Avendo ora a disposizione le conoscenze preliminari necessarie a comprendere l'impiego dei principali strumenti utilizzati per sviluppare Sapienza-DC, è bene effettuare, prima dell'implementazione vera e propria dell'applicativo, alcune dovute considerazioni riguardo la progettazione delle componenti fondamentali che andranno a costituire l'architettura finale del progetto.

Si noti che, come già menzionato, la metodologia di sviluppo agile non prevede obbligatoriamente una raccolta estensiva di requisiti (se non per quanto concerne le principali funzionalità ricercate nell'applicativo) né tantomeno la realizzazione iniziale di diagrammi e modelli tipici dell'approccio a cascata. Difatti, l'assenza di requisiti altamente specifici richiede uno sviluppo dinamico e soggetto a continui cambiamenti, motivo per cui i diagrammi che proporremo di seguito rispecchiano il risultato di quanto ottenuto a fine progetto.   

\section{Diagramma delle classi UML}
Un diagramma delle classi UML (\textit{Unified Modeling Language}) è un diagramma strutturale utilizzato per definire le principali classi, i loro attributi e le relazioni che legano queste ultime. Durante l'esecuzione del software una classe può essere istanziata sotto forma di uno o più oggetti, che definiscono il livello estensionale di quest'ultima. 

Nella programmazione orientata agli oggetti, il diagramma delle classi diviene uno strumento particolarmente utile per disporre di una visione complessiva riguardo il funzionamento del sistema nonché delle relazioni che sussistono tra le varie componenti che lo costituiscono.

A puro scopo nozionistico, si noti che ciascuna classe nel diagramma è contraddistinta da un nome, una serie di attributi locali, ciascuno con il proprio tipo, e da una serie di operazioni associate ad essa.

Per evitare di appesantire il diagramma UML che presenteremo di seguito, si è deciso di non includere nelle specifiche delle operazioni di ciascuna classe i tipi dei parametri in input, così come si è preferito sostituire con la dicitura \texttt{(...)} i parametri accettati dalle operazioni, qualora ve ne fossero in numero elevato.

\begin{sidewaysfigure}[htbp]
    \centering
    \includegraphics[width=1.125\textheight]{diagramma_UML_classi_sdc.png}
    \caption{Diagramma delle classi UML del progetto.}
    \label{fig:diagramma_classi_UML_sdc}
\end{sidewaysfigure}

\clearpage 


\section{Modello ER}
Il modello Entità-Relazione (ER) è un modello di dati impiegato per la progettazione concettuale di basi di dati. La rappresentazione di tali schemi concettuali avviene tipicamente in modo grafico, mediante diagrammi Entità-Relazione. Ogni schema descrive l'aspetto intensionale dei dati; a uno stesso schema possono corrispondere, nel tempo, più istanze, le quali ne costituiscono invece l'aspetto estensionale.

Analogamente a quanto illustrato per i diagrammi delle classi UML, anche il modello ER introduce costrutti simili. Tra questi figurano le entità (analogo delle classi UML), gli attributi di entità, le relazioni (o associazioni) a cui ciascuna entità coinvolta partecipa con il rispettivo ruolo, gli attributi di relazione, le gerarchie di generalizzazione e specializzazione (IS-A), nonché alcuni vincoli di integrità caratteristici del modello Entità-Relazione.

Come per il diagramma delle classi UML, anche per lo schema concettuale del modello ER, ottenuto mediante la corrispettiva progettazione concettuale, verrà fornito direttamente nella sua versione finale, ossia quanto ottenuto al termine dell'implementazione di Sapienza-DC.

\subsection{Progettazione concettuale}
La progettazione concettuale prevede la traduzione dei requisiti forniti nel relativo schema concettuale del modello ER. Nel nostro caso, poiché l'applicazione non deve gestire quantità ingenti di dati, ma solo tenere traccia delle conversazioni e dei messaggi inviati dall'utente, nonché garantire il salvataggio dei dati estratti dalle pagine web (ognuna avente un proprio URL), ai fini di \textit{caching}\footnote{caching: processo di archiviazione temporanea di copie di dati, effettuato con lo scopo di garantire latenza minore in fase di recupero delle informazioni.}, lo schema concettuale risultante è piuttosto semplice.

\begin{figure}[H]
    \centering
    \includegraphics[width=1.025\textwidth]{schema_concettuale_SDC.png}
    \caption{Schema concettuale del modello ER.}
    \label{fig:diagramma_schema_concettuale}
\end{figure}


\subsection{Progettazione logica}
La progettazione logica permette di tradurre lo schema concettuale ottenuto al termine della fase di progettazione concettuale in uno schema relazionale per la nostra base di dati. 

Un accorgimento di vitale importanza da tenere presente in questa fase della progettazione è il seguente: si vuole evitare, ad ogni costo, la perdita di informazioni nel passaggio da uno schema all'altro, affinché anche lo schema relazionale continui a rispettare i requisiti imposti inizialmente. 

\subsubsection{Ristrutturazione dello schema concettuale}
La prima fase della progettazione logica prevede la ristrutturazione dello schema concettuale precedentemente ottenuto. Essa si articola nei seguenti sei passaggi:

\begin{enumerate}
    \item analisi delle ridondanze;
    \item eliminazione degli attributi multivalore;
    \item eliminazione degli attributi composti;
    \item eliminazione delle IS-A e di eventuali generalizzazioni;
    \item scelta degli identificatori principali di entità e relazioni;
    \item specifica degli ulteriori vincoli esterni.
\end{enumerate}

Giacché lo schema concettuale da ristrutturare è elementare, l'unico passo applicabile è il quinto. In particolare, si sceglie come identificatore principale dell'entità \texttt{ParsedURL} l'attributo \textit{id}. Le rimanenti scelte sono immediate e non vengono dunque esplicitate. Lo schema ristrutturato appare pertanto fondamentalmente invariato, così come i vincoli esterni definiti in precedenza.      

\begin{figure}[H]
    \centering
    \includegraphics[width=1.025\textwidth]{schema_concettuale_ristrutturato_SDC.png}
    \caption{Schema concettuale ristrutturato.}
    \label{fig:diagramma_schema_concettuale_ristrutturato}
\end{figure}

\subsubsection{Traduzione diretta}
Lo scopo della traduzione diretta è quello di trasformare lo schema concettuale ristrutturato in uno schema relazionale. Ciascuna relazione viene rappresentata secondo la seguente notazione:
\[\mathscr{R}(\underline{\mathscr{A}_1}, \underline{\mathscr{A}_2}, \mathscr{A}_3, \mathscr{A}^{*}_4,\, \dots \,, \mathscr{A}_n)\]
dove gli attributi sottolineati indicano la chiave primaria della relazione; mentre quelli indicati con l'asterisco sono attributi per cui è ammesso il valore nullo, contrariamente agli attributi non contrassegnati mediante quest'ultimo.

Altre informazioni essenziali come eventuali vincoli di \textit{foreign key}\footnote{foreign key: è un vincolo di integrità referenziale definito fra una sequenza non vuota $\mathcal{X}$ di $n$ attributi di
una relazione $\mathscr{R}_1$ ed una sequenza $\mathcal{Y}$ di $n$ attributi che formano
una chiave di una relazione $\mathscr{R}_2$; è un'asserzione che impone
che per ogni istanza $\mathscr{B}$ della base di dati ogni combinazione di
valori su $\mathcal{X}$ che non contengono NULL e che sono presenti
nella istanza di $\mathscr{R}_1$ in $\mathscr{B}$ compaia come combinazione di valori
su $\mathcal{Y}$ nella istanza di $\mathscr{R}_2$ in $\mathscr{B}$.}, inclusioni, vincoli di chiave oppure ulteriori vincoli di tupla (o inter-tupla) sono da definire al di sotto della relazione stessa.

Una volta terminata la traduzione diretta, otteniamo il seguente schema $\mathcal{S}$ per la nostra base di dati $\mathscr{B}$:\\

\indent
\textbf{Session}(\underline{session\_id}, title*, last\_domain*, created\_at)\\
\indent\indent vincolo di tupla: title is NULL $\Rightarrow $ title = ``New Chat''\\
\indent\indent vincolo di tupla: last\_domain is NULL $\Rightarrow$ last\_domain = ``''\\
\indent\indent vincolo di tupla: last\_domain is not NULL $\Rightarrow$ last\_domain $\in$ \{``www.marvel.com'',\\ 
\indent\indent\indent\indent\indent\indent\,\,\,\, ``www.ipsos.com'', ``www.raiplaysound.com'', ``it.wikipedia.org'',\\
\indent\indent\indent\indent\indent\indent\,\,\,\, ``*''\}\\
\indent
\textbf{Message}(\underline{id}, role, content, created\_at)\\
\indent\indent foreign key: Message[id] $\subseteq$ Contains[message]\\
\indent\indent vincolo di tupla: role $\in$ \{``user'', ``assistant''\}\\
\indent
\textbf{ParsedURL}(\underline{id}, url, parsed\_text*, created\_at)\\
\indent\indent chiave: url\\
\indent\indent vincolo esterno: parsed\_text is not NULL $\Rightarrow$ parsed\_text is in MD format\\
\indent\indent vincolo esterno: url is a valid URL\\ 
\indent
\textbf{Contains}(\underline{message}, session)\\
\indent\indent foreign key: Contains[message] $\subseteq$ Message[id]\\
\indent\indent foreign key: Contains[session] $\subseteq$ Session[session\_id]\\
\indent\indent vincolo esterno: $(m, s) \in \text{Contains} \Rightarrow \text{timestamp}(m) > \text{timestamp}(s)$\\

Si noti infine che si è scelto di assegnare, per comodità, opportuni valori di default agli attributi \textit{title} e \textit{last\_domain} della tabella \texttt{Session}.

\subsubsection{Ristrutturazione dello schema logico}
L'ultima fase della progettazione logica prevede una ristrutturazione dello schema relazionale per motivi di efficienza. Ai fini della nostra applicazione, è di interesse conoscere, ogni qualvolta si accede a un messaggio $m$, la sessione $s$ a cui esso appartiene.

Notando a questo punto che le tabelle \texttt{Message} e \texttt{Contains} sono fortemente accoppiate, si predispone l'accorpamento di \texttt{Contains} a \texttt{Message}, ottenendo il seguente schema relazionale finale per la nostra base di dati:\\

\indent
\textbf{Session}(\underline{session\_id}, title*, last\_domain*, created\_at)\\
\indent\indent vincolo di tupla: title is NULL $\Rightarrow $ title = ``New Chat''\\
\indent\indent vincolo di tupla: last\_domain is NULL $\Rightarrow$ last\_domain = ``''\\
\indent\indent vincolo di tupla: last\_domain is not NULL $\Rightarrow$ last\_domain $\in$ \{``www.marvel.com'',\\ 
\indent\indent\indent\indent\indent\indent\,\,\,\, ``www.ipsos.com'', ``www.raiplaysound.com'', ``it.wikipedia.org'',\\
\indent\indent\indent\indent\indent\indent\,\,\,\, ``*''\}\\
\indent
\textbf{Message}(\underline{id}, session, role, content, created\_at)\\
\indent\indent foreign key: Message[session] $\subseteq$ Session[session\_id]\\
\indent\indent vincolo di tupla: role $\in$ \{``user'', ``assistant''\}\\
\indent\indent vincolo esterno: $m \in \text{Message} \Rightarrow \text{m.created\_at} > \text{timestamp}(\text{m.session})$\\
\indent
\textbf{ParsedURL}(\underline{id}, url, parsed\_text*, created\_at)\\
\indent\indent chiave: url\\
\indent\indent vincolo esterno: parsed\_text is not NULL $\Rightarrow$ parsed\_text is in MD format\\
\indent\indent vincolo esterno: url is a valid URL\\ 

La progettazione della nostra base di dati è a questo punto ultimata. Per completezza, si fornisce di seguito il codice SQL per la definizione delle tabelle del database definitivo.


\begin{lstlisting}[language=SQL, caption={Codice SQL per la creazione delle tabelle della base di dati.}, label={lst:init_db}]
CREATE TABLE IF NOT EXISTS sessions (
    session_id VARCHAR(36) PRIMARY KEY,
    title VARCHAR(255) DEFAULT "New Chat",
    last_domain VARCHAR(255) DEFAULT "",
    created_at DATETIME DEFAULT CURRENT_TIMESTAMP
);

CREATE TABLE IF NOT EXISTS messages (
    id INT AUTO_INCREMENT PRIMARY KEY,
    session_id VARCHAR(36) NOT NULL,
    role VARCHAR(20) NOT NULL,
    content TEXT NOT NULL,
    created_at DATETIME DEFAULT CURRENT_TIMESTAMP,
    FOREIGN KEY (session_id) REFERENCES sessions(session_id) ON 
    DELETE CASCADE
);

CREATE TABLE IF NOT EXISTS parsed_urls (
    id INT AUTO_INCREMENT PRIMARY KEY,
    url TEXT NOT NULL UNIQUE,
    parsed_text LONGTEXT,
    created_at DATETIME DEFAULT CURRENT_TIMESTAMP
);

\end{lstlisting}

Si noti infine l'aggiunta della clausola \texttt{ON DELETE CASCADE}, la quale predispone l'eliminazione automatica e coerente di tutti i messaggi associati a una determinata sessione, nel caso in cui quest'ultima dovesse essere eliminata. Questa scelta permette di mantenere, in ogni istante, l'integrità referenziale del database.   


